% Exact Gamma-Fiber Local Density and the Kummer-Character Remainder
% J. Wilder, 2026.

\documentclass[12pt]{amsart}
\usepackage[utf8]{inputenc}
\usepackage[T1]{fontenc}
\usepackage{textcomp}
\usepackage{amsmath, amssymb, amsthm}
\usepackage[margin=1in]{geometry}

\newtheorem{lemma}{Lemma}
\newtheorem{theorem}{Theorem}
\newtheorem{proposition}{Proposition}
\newtheorem{corollary}{Corollary}
\theoremstyle{definition}
\newtheorem{definition}{Definition}
\theoremstyle{remark}
\newtheorem{remark}{Remark}

\title{Exact $\Gamma$-Fiber Local Density and the Kummer-Character Remainder for $m\cdot 2^k 3^\ell+1$}
\author{Jared Wilder}
\date{2026-05-13}

\begin{document}
\maketitle

\begin{abstract}
We replace the slice-jet Stepanov $|B_q(c)| \ll \sqrt{h_q}\log q$ bound with the exact
group-theoretic formula $|B_q(c)| = \gcd(a_q, b_q) \cdot \mathbf{1}_{\Gamma_q}(c)$,
where $\Gamma_q = \langle 2, 3 \rangle \subset \mathbb{F}_q^*$. By character orthogonality
the local divisibility density splits as $\rho_m(q) = 1/(q-1) + \Delta_m(q)$, with
$\Delta_m(q)$ a Kummer-character discrepancy that admits an explicit closed form,
a Pappalardi-controlled pointwise bound, and empirical verification of cancellation
on average over $q$. The remaining analytic problem is recast as a
Bombieri-Vinogradov-style estimate for the Kummer-discrepancy sum.
\end{abstract}

%% ===================================================================
\section{Notation}
%% ===================================================================

\begin{definition}
Let $q \nmid 6m$ be prime. Set
$$
H_2(q) = \langle 2 \rangle \subset \mathbb{F}_q^*, \quad H_3(q) = \langle 3 \rangle \subset \mathbb{F}_q^*,
$$
$$
a_q = \mathrm{ord}_q(2), \quad b_q = \mathrm{ord}_q(3),
$$
$$
\Gamma_q = \langle 2, 3 \rangle = H_2(q) H_3(q) \subset \mathbb{F}_q^*,
$$
$$
B_q(c) = \{(x, y) \in H_2(q) \times H_3(q) : x y \equiv c \pmod q\}\ \text{for}\ c \in \mathbb{F}_q^*.
$$
\end{definition}

%% ===================================================================
\section{Exact $\Gamma$-fiber formula}
%% ===================================================================

\begin{lemma}[Exact slice cardinality]\label{lem:R682}
For every prime $q$ with $q \nmid 6m$ and $c \in \mathbb{F}_q^*$,
$$
|B_q(c)|
\;=\;
\gcd(a_q, b_q) \cdot \mathbf{1}_{\Gamma_q}(c).
$$
\end{lemma}

\begin{proof}
Consider the multiplication homomorphism $\mu : H_2(q) \times H_3(q) \to \mathbb{F}_q^*$,
$(x, y) \mapsto xy$. Its image is $\Gamma_q = H_2(q) H_3(q)$. If $c \notin \Gamma_q$,
$|B_q(c)| = |\mu^{-1}(c)| = 0$. If $c \in \Gamma_q$, every nonempty fibre of a group
homomorphism has size $|\ker \mu|$. Now
$\ker \mu = \{(t, t^{-1}) : t \in H_2(q) \cap H_3(q)\}$, and since $\mathbb{F}_q^*$ is
cyclic, $|H_2(q) \cap H_3(q)| = \gcd(a_q, b_q)$.
\end{proof}

\begin{corollary}\label{cor:Gamma-size}
$|\Gamma_q| = a_q b_q / \gcd(a_q, b_q) = \mathrm{lcm}(a_q, b_q)$.
\end{corollary}

\begin{remark}[Empirical verification]
The formula was verified for all $q \in [5, 200]$ prime, $m \in \{1, 5, 7, 11, 13, 17\}$:
$259/259$ (q, m) pairs satisfy the identity exactly, by direct enumeration of
$\{(k, \ell) \in [a_q] \times [b_q] : 2^k 3^\ell \equiv -m^{-1}\}$ vs.\
$\gcd(a_q, b_q) \cdot \mathbf{1}_{\Gamma_q}(-m^{-1})$.
\hfill(Numerical verification code accompanies this note.)
\end{remark}

%% ===================================================================
\section{Box-period local count}
%% ===================================================================

\begin{lemma}[Exact prime-modulus box count]\label{lem:R683}
For $K, L \ge 1$ and $c_q = -m^{-1} \bmod q$,
$$
A_m(q; K, L)
\;:=\;
\# \{0 \le k < K,\ 0 \le \ell < L : m 2^k 3^\ell + 1 \equiv 0 \pmod q\}
\;=\;
\frac{KL}{a_q b_q} |B_q(c_q)|
\;+\;
O(a_q L + b_q K + a_q b_q).
$$
Using Lemma~\ref{lem:R682},
$$
A_m(q; K, L)
\;=\;
\frac{KL}{\mathrm{lcm}(a_q, b_q)} \cdot \mathbf{1}_{\Gamma_q}(-m^{-1})
\;+\;
O(a_q L + b_q K + a_q b_q).
$$
\end{lemma}

\begin{proof}
The pair $(2^k \bmod q, 3^\ell \bmod q)$ is periodic in $(k, \ell)$ modulo $(a_q, b_q)$.
Decompose $[K] \times [L]$ into complete fundamental rectangles of size $a_q \times b_q$
plus boundary strips of size $\le a_q L + b_q K + a_q b_q$.
Each complete rectangle contains exactly $|B_q(c_q)|$ solutions.
\end{proof}

%% ===================================================================
\section{Character expansion}
%% ===================================================================

\begin{lemma}[Kummer-character identity]\label{lem:R684}
Let $\Gamma_q^\perp = \{\chi \bmod q : \chi(2) = \chi(3) = 1\}$. Then
$$
\frac{\mathbf{1}_{\Gamma_q}(c)}{|\Gamma_q|}
\;=\;
\frac{1}{q - 1}
\sum_{\chi \in \Gamma_q^\perp}
\chi(c).
$$
Setting $c = -m^{-1}$, the local divisibility density splits as
$$
\rho_m(q) := \frac{|B_q(-m^{-1})|}{a_q b_q}
= \frac{\mathbf{1}_{\Gamma_q}(-m^{-1})}{|\Gamma_q|}
= \frac{1}{q-1}
+ \Delta_m(q),
$$
where the Kummer-character discrepancy is
$$
\Delta_m(q)
\;=\;
\frac{1}{q-1}
\sum_{\substack{\chi \in \Gamma_q^\perp \\ \chi \ne \chi_0}}
\chi(-m^{-1}).
$$
\end{lemma}

\begin{proof}
For any subgroup $\Gamma \le G$, the annihilator $\Gamma^\perp \le \widehat{G}$ satisfies
$|\Gamma^\perp| = |G|/|\Gamma|$. Character orthogonality on $G/\Gamma$ gives
$\mathbf{1}_\Gamma(c) = (|\Gamma|/|G|) \sum_{\chi \in \Gamma^\perp} \chi(c)$.
With $|G_q| = q - 1$ and dividing by $|\Gamma_q|$ yields the first identity.
The split $\rho_m(q) = 1/(q-1) + \Delta_m(q)$ follows by separating the principal character.
\end{proof}

%% ===================================================================
\section{$\Delta_m(q)$ closed form and Pappalardi bound}
%% ===================================================================

\begin{lemma}[Closed form for $\Delta_m(q)$]\label{lem:Delta-closed}
For every prime $q \nmid 6m$,
$$
\Delta_m(q)
=
\begin{cases}
\dfrac{1}{|\Gamma_q|} - \dfrac{1}{q-1}
& \text{if } -m^{-1} \in \Gamma_q,
\\[1.2em]
-\dfrac{1}{q-1}
& \text{if } -m^{-1} \notin \Gamma_q.
\end{cases}
$$
In particular,
$$
|\Delta_m(q)|
\;\le\;
\frac{1}{|\Gamma_q|} - \frac{1}{q-1}
\;\le\;
\frac{1}{|\Gamma_q|}.
$$
\end{lemma}

\begin{proof}
By Lemma~\ref{lem:R684}, $\rho_m(q) = \mathbf{1}_{\Gamma_q}(-m^{-1}) / |\Gamma_q|$.
\begin{itemize}
\item If $-m^{-1} \in \Gamma_q$, then $\rho_m(q) = 1/|\Gamma_q|$, so
$\Delta_m(q) = 1/|\Gamma_q| - 1/(q-1) \ge 0$ since $|\Gamma_q| \le q - 1$.
\item If $-m^{-1} \notin \Gamma_q$, then $\rho_m(q) = 0$, so $\Delta_m(q) = -1/(q-1)$.
\end{itemize}
The pointwise bound is immediate.
\end{proof}

\begin{lemma}[Pappalardi-controlled bound]\label{lem:Delta-Pap}
Let $\mathcal{E}^{\mathrm{Pap}}(Q) = \{q \le Q : |\Gamma_q| \le q^{1/2}\}$.
By Pappalardi (1996, J.~Number Theory 57, Thm.~1, instantiated at rank $r = 2$, $\alpha = 1/2$),
$$
|\mathcal{E}^{\mathrm{Pap}}(Q)| \;\ll\; Q^{1 - \eta + \varepsilon},
\quad
\eta = 1/12,
\quad
\text{unconditionally, no GRH.}
$$
Hence, for every prime $q \le Q$ with $q \notin \mathcal{E}^{\mathrm{Pap}}(Q)$:
$$
|\Delta_m(q)| \;\le\; \frac{1}{|\Gamma_q|} \;<\; q^{-1/2},
$$
uniformly in $m$ coprime to $6$.
\end{lemma}

\begin{proof}
For $q \notin \mathcal{E}^{\mathrm{Pap}}(Q)$, $|\Gamma_q| > q^{1/2}$ by definition.
Combine with Lemma~\ref{lem:Delta-closed}.
\end{proof}

\begin{lemma}[Average over $m$ vanishes]\label{lem:Delta-avg}
For every fixed prime $q \nmid 6$,
$$
\frac{1}{q - 1}
\sum_{\substack{m \in \mathbb{F}_q^* \\ \gcd(m, 6) = 1}}
\Delta_m(q)
\;=\;
0.
$$
\end{lemma}

\begin{proof}
By Lemma~\ref{lem:Delta-closed},
$$
\sum_m \Delta_m(q) = (\text{$\#\{m : -m^{-1} \in \Gamma_q\}$}) \left(\tfrac{1}{|\Gamma_q|} - \tfrac{1}{q-1}\right) + (\text{$\#\{m : -m^{-1} \notin \Gamma_q\}$}) \left(- \tfrac{1}{q-1}\right).
$$
The map $m \mapsto -m^{-1}$ is a bijection on $\mathbb{F}_q^*$, so
$\#\{m : -m^{-1} \in \Gamma_q\} = |\Gamma_q|$ and
$\#\{m : -m^{-1} \notin \Gamma_q\} = q - 1 - |\Gamma_q|$. Substituting gives
$$
\sum_m \Delta_m(q) = |\Gamma_q| \left(\tfrac{1}{|\Gamma_q|} - \tfrac{1}{q-1}\right) - (q - 1 - |\Gamma_q|) \tfrac{1}{q-1}
= 1 - \tfrac{|\Gamma_q|}{q-1} - 1 + \tfrac{|\Gamma_q|}{q-1} = 0.
\qedhere
$$
\end{proof}

\begin{remark}[Empirical cancellation]
For $q \le 200$ and $m \in \{1, 5, 7, 11, 13, 17\}$, the partial sums
$\sum_{q \le Q} \Delta_m(q)$ remain bounded (RMS $|\Delta_m(q)| \approx 0.01$ at $Q = 200$,
total $\sum_q |\Delta_m(q)| \le 0.07$ across $m$). The stable boundedness as $Q$ grows is
consistent with genuine BV-style cancellation, not merely the Pappalardi triangle bound.
\end{remark}

%% ===================================================================
\section{Distribution formula}
%% ===================================================================

\begin{lemma}[Corrected distribution formula]\label{lem:R685}
$$
A_m(q; K, L)
\;=\;
\frac{KL}{q - 1}
\;+\;
KL \cdot \Delta_m(q)
\;+\;
O(a_q L + b_q K + a_q b_q).
$$
\end{lemma}

\begin{proof}
Combine Lemmas \ref{lem:R683} and \ref{lem:R684}.
\end{proof}

%% ===================================================================
\section{Squarefree modulus}
%% ===================================================================

\begin{lemma}[Squarefree modulus distribution]\label{lem:R686}
Let $d$ squarefree, $\gcd(d, 6m) = 1$. Set $\Gamma_d = \langle 2, 3 \rangle \subset (\mathbb{Z}/d\mathbb{Z})^*$,
$\Gamma_d^\perp = \{\chi \bmod d : \chi(2) = \chi(3) = 1\}$. Then
$$
\rho_m(d)
\;=\;
\frac{\mathbf{1}_{\Gamma_d}(-m^{-1})}{|\Gamma_d|}
\;=\;
\frac{1}{\varphi(d)}
\;+\;
\Delta_m(d),
\quad
\Delta_m(d) = \frac{1}{\varphi(d)}
\sum_{\substack{\chi \bmod d \\ \chi \ne \chi_0 \\ \chi(2) = \chi(3) = 1}}
\chi(-m^{-1}).
$$
$$
A_m(d; K, L)
\;=\;
\frac{KL}{\varphi(d)}
\;+\;
KL \cdot \Delta_m(d)
\;+\;
O(a_d L + b_d K + a_d b_d).
$$
Note: $\rho_m(d)$ is generally not multiplicative in $d$ because the character constraint
$\chi(2) = \chi(3) = 1$ is imposed at the level of $d$.
\end{lemma}

%% ===================================================================
\section{Exact Gate 0 theorem}
%% ===================================================================

\begin{theorem}[Bombieri-Vinogradov for Kummer-discrepancy]\label{thm:BV-Kummer}
The level-of-distribution problem for $m 2^k 3^\ell + 1$ reduces to the
Kummer-discrepancy cancellation:
$$
\sum_{d \le D} \lambda_d \Delta_m(d)
\;\ll\;
(\log X)^{-A},
$$
for admissible sieve weights $\lambda_d$ (i.e., $|\lambda_d| \le \tau(d)^{O(1)}$),
uniformly in $m$ coprime to 6.

Conditional on the above, the orbit count satisfies the Bombieri-Vinogradov-type bound
$$
\sum_{d \le D} \lambda_d \left( A_m(d; K, L) - \frac{KL}{\varphi(d)} \right)
\;\ll\;
KL (\log X)^{-A}
\;+\;
O\!\left( \sum_{d \le D} |\lambda_d| (a_d L + b_d K + a_d b_d) \right).
$$
\end{theorem}

\begin{proof}
By Lemma~\ref{lem:R686},
$A_m(d; K, L) - KL/\varphi(d) = KL \Delta_m(d) + O(a_d L + b_d K + a_d b_d)$. Multiply by $\lambda_d$ and sum.
\end{proof}

%% ===================================================================
\section{The Cor 4.1\textquotedblright\ chain (with explicit gates)}
%% ===================================================================

The chain to Corollary 4.1$''$ (unconditional NEG-203) now reads:

\begin{enumerate}
\item Lemma~\ref{lem:R682} (exact fiber): exact group-theoretic identity.
\item Lemma~\ref{lem:R683} (box count): exact box-period count.
\item Lemma~\ref{lem:R684} (character expansion): $\rho_m(q) = 1/(q-1) + \Delta_m(q)$.
\item Lemma~\ref{lem:Delta-closed} ($\Delta$-closed form): pointwise formula for $\Delta_m(q)$.
\item Lemma~\ref{lem:Delta-Pap} (Pappalardi-controlled bound): $|\Delta_m(q)| < q^{-1/2}$ on density-one primes.
\item Lemma~\ref{lem:Delta-avg} ($m$-average vanishes): $\mathbb{E}_m \Delta_m(q) = 0$.
\item Lemma~\ref{lem:R685} (distribution): $A_m(q; K, L) = KL/(q-1) + KL \Delta_m(q) + O(\cdots)$.
\item Lemma~\ref{lem:R686} (squarefree): same for $d$ squarefree.
\item Theorem~\ref{thm:BV-Kummer} (Gate 0): BV-style cancellation of $\sum_d \lambda_d \Delta_m(d)$ \(\Longrightarrow\) orbit level of distribution.
\item Pappalardi $\eta = 1/12$ + Heath-Brown identity $K = 4$ + Type I (Lemma~\ref{lem:R686} + Pappalardi) + Type II (Theorem~\ref{thm:BV-Kummer} + FI ASP threshold) \(\Longrightarrow\) $\sum \Lambda(m 2^k 3^\ell + 1) \to \infty$ \(\Longrightarrow\) $A(m)$ infinite.
\end{enumerate}

The hard remaining gate is Theorem~\ref{thm:BV-Kummer}. The Pappalardi-controlled pointwise
bound (Lemma~\ref{lem:Delta-Pap}) gives only the triangle bound
$\sum_{d \le D} |\Delta_m(d)| \ll D^{1/2} + D^{1-\eta+\varepsilon}$, which is too weak for BV.
Genuine cancellation comes from oscillation in $\chi(-m^{-1})$ across the moduli $d$, which
is the Linnik-Burgess-style character sum problem.

%% ===================================================================
\section{Empirical anchor}
%% ===================================================================

\begin{itemize}
\item Lemma~\ref{lem:R682} verified for all 259 (q, m) pairs with $q \le 200$, $m \in \{1, 5, 7, 11, 13, 17\}$.
\item $\Delta_m(q)$ values computed via the closed-form Lemma~\ref{lem:Delta-closed}.
\item Partial sums $\sum_{q \le Q} \Delta_m(q)$ stable at $\le 0.07$ across $Q \in \{50, 100, 150, 200\}$ for all tested $m$.
\item RMS $|\Delta_m(q)|$ decreases from $\approx 0.014$ at $Q = 50$ to $\approx 0.008$ at $Q = 200$, consistent with $1/\sqrt{Q}$ scaling.
\end{itemize}

%% ===================================================================
\section{Conclusion}
%% ===================================================================

A direct Stepanov auxiliary-polynomial approach to bounding $|B_q(c)|$ is treated in a companion note
and gives the unconditional upper bound $|B_q(c)| \ll \sqrt{h_q}\log q$. In the present setting we
have the sharper exact group-theoretic identity above, and the remaining analytic problem is the
Bombieri--Vinogradov-style cancellation of the Kummer-discrepancy sum $\sum_d \lambda_d \Delta_m(d)$
(Theorem~\ref{thm:BV-Kummer}).

Lemmas~\ref{lem:Delta-closed}, \ref{lem:Delta-Pap}, \ref{lem:Delta-avg}
together supply the closed form, Pappalardi-controlled pointwise bound, and average-zero property of
$\Delta_m(q)$. Direct numerical evaluation confirms the identity and exhibits BV-cancellation patterns
on $q \le 200$.

%% ===================================================================
\section*{References}
%% ===================================================================

\begin{itemize}
\item[{[P96]}] F.~Pappalardi, ``On the order of finitely generated subgroups of $\mathbb{Q}^*$ $\pmod p$ and divisors of $p - 1$,'' J.~Number Theory 57 (1996), 207--222. Thm.~1 instantiated at rank $r = 2$, $\alpha = 1/2$: $\eta = 1/12$.
\item[{[HB96]}] D.R.~Heath-Brown, ``A new $k$-th power identity,'' Proc.~London Math.~Soc.~(3) 73 (1996), 241--301.
\item[{[FI98]}] J.~Friedlander and H.~Iwaniec, ``The polynomial $X^2 + Y^4$ captures its primes,'' Annals of Math.~148 (1998), 1041--1065.
\item[{[BV65]}] E.~Bombieri, ``On the large sieve,'' Mathematika 12 (1965), 201--225.
\item[{[L80]}] H.~Iwaniec, ``A new form of the error term in the linear sieve,'' Acta Arithmetica 37 (1980), 307--320.
\end{itemize}

\end{document}
